\chapter{検証結果}\label{chap:5}

本章では、全54条件を対象とした機能検証の結果、および追加検証・異常系テストの結果を報告する。
続いて、代表入力に対する1コア実行と2コア実行の総サイクル数を比較し、要素数や方式による傾向を示す。

\section{出力と途中状態の検査結果}\label{sec:5-1}

18種類の入力ケースに対して3つの実行方式を適用した計54条件についてシミュレーションを実施した結果、すべての条件で最終配列が期待値と完全に一致し、\ref{sec:3-3}節で述べた実行途中の各種検査を通過した。
各条件の実行ログにおけるPASS出力および終了コード0の確認結果は、総サイクル数とあわせて付録\ref{chap:B}にまとめる。
また、指定10要素をdual方式で実行した際の詳細な波形は、前章の図\ref{fig:3}〜\ref{fig:5}に示したとおりである。

主検証の54条件に加えて、補助的な動作確認も実施した。
具体的には、指定10要素および乱数64要素を各3方式で実行した計6件の波形デモと、コアID読み出し、既存のマルチコア基本動作、ROM読み出し時の競合、およびメモリ配置を対象とした計4件の単体検証を実施し、いずれも正常に通過することを確認した。
これらは主検証の54条件とは独立した実行記録であり、波形デモの入力ケースは主検証と共通である。

さらに、テストベンチの期待値照合が異常を確実に検出できることを確認するため、意図的に誤った期待値を与えた異常系テストを実行した。
このテストでは、array[0]の期待値0xFFFFFFFFに対して実際の出力値が0x00000000である不一致を検出し、終了コード1で実行を停止した。
単一の異常系テストのみで検査器の全エラー経路を網羅できるわけではないが、最終配列の不一致検出機能が意図どおりに働くことを確認できた。

\section{1コア・2コア実行の総サイクル数比較}\label{sec:5-2}

表\ref{tab:5}に、代表入力における各方式の総サイクル数およびserial\_splitとdualのサイクル比を示す。
主比較である1コア逐次分割方式（serial\_split）と2コア並列方式（dual）を比べると、指定10要素では442サイクルから366サイクルへ短縮された（比1.21）。
固定乱数の4ケースでは、サイクル比が7要素の1.08から64要素の1.56へと拡大しており、要素数が増えるにつれて並列実行の短縮効果が顕著になる傾向が確認された。
ただし、各要素数につき1種類の固定配列のみを評価しているため、入力の並び順による変動と要素数による効果は厳密には分離されていない。

\begin{table}[H]
\centering
\caption{代表入力における各方式の総サイクル数}
\label{tab:5}
\footnotesize
\setlength{\tabcolsep}{3pt}
\renewcommand{\arraystretch}{1.25}
\begin{tabular}{>{\raggedright\arraybackslash}p{\dimexpr0.150218\linewidth-2\tabcolsep\relax}>{\raggedright\arraybackslash}p{\dimexpr0.216313\linewidth-2\tabcolsep\relax}>{\raggedright\arraybackslash}p{\dimexpr0.170247\linewidth-2\tabcolsep\relax}>{\raggedright\arraybackslash}p{\dimexpr0.190276\linewidth-2\tabcolsep\relax}>{\raggedright\arraybackslash}p{\dimexpr0.272947\linewidth-2\tabcolsep\relax}}\toprule
\textbf{入力} & \textbf{serial\_split} & \textbf{dual} & \textbf{split/dual} & \textbf{参考 serial\_full}\\\midrule
指定10 & 442 & 366 & 1.21 & 441\\
乱数7 & 281 & 260 & 1.08 & 230\\
乱数17 & 923 & 733 & 1.26 & 1,164\\
乱数32 & 2,513 & 1,772 & 1.42 & 3,763\\
乱数64 & 8,614 & 5,527 & 1.56 & 14,803\\
\bottomrule\end{tabular}
\end{table}

一方、参考条件であるserial\_fullとの比較では、要素数が小さい乱数7においてserial\_fullの230サイクルに対しdualが260サイクルと30サイクル上回った。
これは、要素数が少ない領域では分割とマージおよび同期の固定的なオーバーヘッドがバブルソート自体の処理量に対して大きく影響するためである。
逆に17要素以上では配列分割による比較回数の削減効果が上回り、dualとserial\_splitのいずれもserial\_fullより少ないサイクル数となった。
もっとも、serial\_fullはマージを伴わない点でアルゴリズム構成が異なるため、この差を単純にコア数のみの効果として解釈すべきではない。
